\documentclass[11pt]{article}

\usepackage[utf8]{inputenc}
\usepackage[T1]{fontenc}
\usepackage{lmodern}
\usepackage[margin=1in]{geometry}
\usepackage{amsmath,amssymb}
\usepackage{booktabs}
\usepackage{array}
\usepackage{microtype}
\usepackage{xcolor}
\usepackage[hidelinks]{hyperref}
\usepackage[capitalize,noabbrev]{cleveref}

\newcommand{\bt}[1]{\texttt{#1}}
\newcommand{\TODO}[1]{{\color{red}\textbf{[TODO: #1]}}}

\title{\bfseries One Program, Many Circuits:\\
Backend-Aware Circuit Optimization for Private AI}

\author{(author names elided for review)}
\date{}

\begin{document}
\maketitle

\begin{abstract}
Privacy-preserving machine learning is written by experts, by hand, against a
chosen cryptographic backend. We argue it should be compiled, and that a
compiler can do better than the expert --- not because it is cleverer, but
because batching independent work is a property of a schedule rather than of
an author's attention. We present Baru, a language whose programs compile to
a single circuit IR and are then optimized \emph{differently depending on
which privacy backend they target}. On a neural network with secret weights,
compiled code issues $0.38\times$ the interactive protocol invocations and
$0.49\times$ the messages of a hand-written, hand-batched expert baseline,
while computing a bit-identical result; the source says nothing about
batching. We report an ablation in which each compiler pass acts on a
different axis --- scheduling collapses rounds $285\times$, common
subexpression elimination collapses circuit size $68\times$ and bandwidth
$30\times$, and rebalancing dependent reductions turns a fold's $2n-1$ rounds
into $2\lceil\log_2 n\rceil+1$. Finally we show the property that motivates
the design: the \emph{same} rewrite is a $2.24\times$ win under one protocol
and worth exactly nothing under another, so a single optimized circuit cannot
serve both. Optimizing one program per backend yields up to three distinct
circuits, and running a circuit under the backend it was not optimized for
costs up to $2.33\times$.
\end{abstract}

\section{Introduction}\label{sec:intro}

Systems for private machine learning --- secure multi-party computation (MPC),
homomorphic encryption (FHE), and hybrids of the two --- are overwhelmingly
written by hand. The reason is not ignorance of compilers; it is that the
performance of these systems depends on details a general-purpose compiler is
not aware of. Under a secret-sharing protocol the dominant online cost is not
work but \emph{communication rounds}, and the difference between a layer that
costs one round and one that costs $n$ is invisible to the source text.

This paper starts from an observation about our own code.

\paragraph{Experts serialize layers, and do not notice.} Our hand-written
private-inference pipeline contained a scheme-switch import that issued three
blocking communication rounds \emph{per activation} --- $16\,224$ sequential
round trips for a single convolution. It accounted for $88\%$ of end-to-end
runtime. It survived review, survived benchmarking, and a draft of a
companion paper had already concluded in print that the bottleneck lay
elsewhere. It was found only by per-phase instrumentation.

That was not an isolated slip. In the same codebase, the composable gadget
library --- the API a user is meant to write against --- issued one comparison
and one multiplication \emph{per element} for ReLU, and one multiplication
\emph{per weight} for a dense layer. Measured over a real transport on a
$64$-wide layer, this cost $5$--$14\times$ the messages and $9$--$45\times$
the wall-clock of the same computation issued through batched primitives, with
byte counts within $5\%$ and outputs bit-identical. The defect is invisible to
a security argument, invisible to a correctness test, and invisible at zero
latency. Only a round count exposes it.

\paragraph{The claim.} A compiler does not have this failure mode. An
as-soon-as-possible schedule over the dependence graph batches
\emph{everything} that is independent, uniformly, without being asked. We show
that compiled code therefore matches and then beats expert hand-written code
on a real network, and --- the part we believe is new --- that the right
circuit \emph{depends on the backend}, so a compiler that optimizes once and
emits one circuit is leaving a measurable amount on the table.

\paragraph{Contributions.}
\begin{enumerate}
\item \textbf{Parity, then advantage (\Cref{sec:eval-network}).} Compiled code
matches hand-batched code message for message at the layer level, and on a
whole network with secret weights issues $0.38\times$ the interactive protocol
invocations and $0.49\times$ the messages of the expert baseline.

\item \textbf{An ablation with separable axes (\Cref{sec:ablation}).} Each
pass contributes on a different metric: scheduling $285\times$ on rounds, CSE
$68\times$ on circuit size and $30\times$ on bandwidth, reduction rebalancing
$1.25\times$ on rounds. None substitutes for another.

\item \textbf{Rebalancing \emph{compound} reductions (\Cref{sec:reassoc}).}
Classical reassociation handles chains of a single associative primitive. The
reductions an AI workload contains are not primitives --- a maximum is
$\mathrm{acc} + (x>\mathrm{acc})\cdot(x-\mathrm{acc})$, four instructions with
two entry wires. We recognise an associative reduction by repeated subgraph
shape, verify associativity by evaluation, and rebalance: $2n-1$ rounds become
$2\lceil\log_2 n\rceil+1$.

\item \textbf{Backend-aware optimization, measured
(\Cref{sec:backend-aware}).} We optimize one program once per backend and
price every resulting circuit under every backend. The diagonal wins every
column; the off-diagonal entries are what a backend-blind compiler loses.
\end{enumerate}

\section{Background and setting}\label{sec:background}

\paragraph{Rounds, not operations.} In secret-sharing MPC, local arithmetic is
free and every multiplication or comparison requires interaction. Because
independent interactive operations can be merged into one round, the cost of a
program is governed by the \emph{depth} of its dependence graph, not its size.
A layer of $512$ independent multiplications costs one round; a chain of $512$
dependent ones costs $512$.

\paragraph{Backends differ in what they charge for.} An arithmetic
secret-sharing protocol pays per round and treats addition as free. A binary
protocol pays per AND gate and has real depth inside an adder. A
constant-round garbled protocol pays for gates but not for depth at all. A
rewrite that trades depth for nothing is therefore a large win for the first,
a moderate one for the second, and worthless for the third --- from one rule.

\paragraph{Threat model and scope.} All results are semi-honest. Protocol
availability follows the deployment (party count, majority assumption); the
comparisons in \Cref{sec:eval} use the protocols named in each table. Our
executed measurements use a three-party honest-majority replicated protocol~\cite{ABY3};
the analytic cost model spans a wider registry and is the basis of
\Cref{sec:backend-aware}. \Cref{sec:limitations} states which numbers are
measured and which are modelled.

\section{Baru}\label{sec:design}

Baru is a small expression language. A program is compiled to a static
single-assignment circuit IR (\bt{FlatOp}) shared by every backend, then
lowered to a per-protocol instruction program, scheduled, and either costed or
executed.

\[
\text{Expr} \;\to\; \text{FlatOp (SSA)} \;\to\; \text{lower} \;\to\;
\text{schedule} \;\to\; \{\text{cost},\ \text{execute}\}
\]

\paragraph{Scheduling is where batching comes from.} The schedule is ASAP over
the dependence DAG: a local instruction inherits the level of its latest
operand; an interactive instruction sits one level later, and every
interactive instruction at a level executes in one round. ASAP is optimal for
depth, and depth is the round count, so this is the minimum-round schedule
\emph{for a fixed lowering}. Crucially, no vectorized IR is required:
independent scalar operations land at the same level and are issued through
one batched primitive.

\paragraph{Optimization is search plus a cost model.} Rewrites are proposed
per site and accepted when the backend's cost model says they pay. The rules
used here are algebraic identities, common subexpression elimination, shift
fusion, strength reduction, classical reassociation, and reduction
rebalancing (\Cref{sec:reassoc}). Rules whose sign depends on the operand are
searched site by site, so a rewrite can be taken at one site and declined at
another within the same circuit.

\subsection{Rebalancing compound reductions}\label{sec:reassoc}

Classical reassociation rebuilds $((a\cdot b)\cdot c)\cdot d$ into
$(a\cdot b)\cdot(c\cdot d)$ for a single associative operator. AI reductions
are not single operators. A maximum-fold step is
\[
\mathrm{acc}' \;=\; \mathrm{acc} + (x > \mathrm{acc})\cdot(x - \mathrm{acc}),
\]
four instructions forming a subgraph with two entry wires. Our rule peels a
\emph{step region} --- the instructions between the accumulator wire and the
step's result --- requiring that every instruction inside is private to the
region and that exactly two wires enter it. Steps are canonicalised to a shape
key; a chain of equal-keyed steps is a reduction, and its leaves are
rebalanced into a tournament preserving left-to-right order, so the rewrite
needs associativity and never commutativity.

\paragraph{Associativity is tested, not assumed.} The search accepts
candidates on \emph{cost}, not on semantics, so a rule cannot rely on it to
reject an unsound rewrite. Before offering the rewrite we evaluate the step
operator on edge values and random triples, checking $f(f(a,b),c) = f(a,f(b,c))$;
a subtraction fold is correctly left alone. The optimizer additionally
differentially checks the whole resulting circuit.

Three implementation subtleties determined whether the rule worked at all, and
we record them because each produced a plausible wrong answer rather than a
failure. (i) The accumulator is read three times \emph{inside its own step},
so ``used once'' is the wrong privacy test; the correct condition is that
every use lies inside the step that consumes it. (ii) A fold's inner roots are
also valid chains, so taking the first match rebalances a prefix and leaves
the tail serial. (iii) A rebalanced tree still \emph{contains} chains --- every
root-to-leaf path is one --- so a fixpoint pass will happily rebuild a
balanced path into a deeper one; the rule must refuse any rewrite that does
not strictly reduce depth.

\section{Evaluation}\label{sec:eval}

Every figure is produced by a command in the artifact, and every route in
every comparison is checked against a plaintext integer reference --- a cost
comparison between routes computing different things is meaningless, and our
harnesses fail rather than report one. Wall-clock is single-threaded on
commodity hardware; ratios transfer, absolute times do not.

\subsection{Scheduling batches without being asked}\label{sec:eval-sched}

\Cref{tab:flat} compiles layer-shaped programs and reports the ASAP round
count. Interactive operations grow by more than two orders of magnitude while
the round count does not move.

\begin{table}[t]
\centering
\caption{Round count is flat in layer width. No vectorized IR, no
annotation.}
\label{tab:flat}
\begin{tabular}{@{}lrrr@{}}
\toprule
program & $n$ & interactive ops & rounds \\
\midrule
\bt{relu-layer} & $1$ & $3$ & $3$ \\
\bt{relu-layer} & $64$ & $129$ & $3$ \\
\bt{relu-layer} & $512$ & $1025$ & $3$ \\
\bt{dot} & $512$ & $513$ & $2$ \\
\bottomrule
\end{tabular}
\end{table}

\subsection{A whole network against an expert}\label{sec:eval-network}

The network is $x(n) \to \mathrm{dense}(h) \to \mathrm{ReLU} \to
\mathrm{dense}(c) \to \max$, with \emph{secret weights} throughout, so every
multiplication is a real interactive operation. Two routes evaluate it over
the same three-party transport:

\begin{itemize}
\item \textbf{gadget} --- hand-written code calling a batched gadget library.
This is the honest expert baseline: those gadgets match hand-tuned code
message for message.
\item \textbf{compiled} --- the same network as a Baru expression, through the
circuit compiler and a level-batched executor. The source is scalar
arithmetic, a mux-style ReLU and a left-fold maximum. It says nothing about
batching.
\end{itemize}

At $n{=}32$, $h{=}16$, $c{=}8$, RTT $20$ ms: the expert route issues $39$
interactive protocol invocations and $95$ messages; the fully optimized
compiled route issues $\mathbf{15}$ and $\mathbf{47}$ --- $0.38\times$ and
$0.49\times$ --- at parity on wall-clock. The advantage \emph{grows with
width}: the compiled route stays at $15$ invocations while the expert route
grows from $31$ at $n{=}16$ to $39$ at $n{=}32$.

\paragraph{Two things we do not claim.} Without reduction rebalancing the
compiled route \emph{loses} ($1.88\times$ wall-clock), because the expert's
maximum is a hand-written tournament; that ablation is reported rather than
hidden. And the wall-clock is parity, not a win: both routes perform the same
number of comparison layers, which dominate latency under a secret-sharing
protocol. The compiler's gain is in multiplications and sharing.

\subsection{What each pass contributes}\label{sec:ablation}

\Cref{tab:ablation} adds one pass per row, starting from an executor written
the obvious way: walk the circuit, call the backend per operation.

\begin{table}[t]
\centering
\caption{Ablation, $n{=}8$, $h{=}4$, $c{=}4$, zero latency. Every row computes
the same value.}
\label{tab:ablation}
\begin{tabular}{@{}lrrrr@{}}
\toprule
route & ops & interactive & messages & bytes \\
\midrule
\bt{gadget} (expert) & --- & $20$ & $53$ & $56$\,k \\
\midrule
\bt{compiled: serial} & $7419$ & $4280$ & $5859$ & $1798$\,k \\
\quad $+$ ASAP batching & $7419$ & $15$ & $47$ & $1682$\,k \\
\quad $+$ CSE & $109$ & $15$ & $47$ & $55.7$\,k \\
\quad $+$ reassociation & $109$ & $\mathbf{12}$ & $\mathbf{37}$ & $55.5$\,k \\
\bottomrule
\end{tabular}
\end{table}

The passes act on \emph{different axes}: scheduling is $285\times$ on
interactive invocations and $1.00\times$ on circuit size; CSE is $68\times$ on
circuit size and $30\times$ on bytes while changing no round count;
reassociation is $1.25\times$ on rounds and nothing else. A single-column
summary would report CSE as a $1.00\times$ no-op, which is why we report the
ladder.

The \bt{serial} row deserves emphasis. It is $285\times$ the interactive
invocations of the \emph{same circuit} under ASAP levelling, on essentially
identical bytes. The round-structure defect that our human experts shipped
(\Cref{sec:intro}) reappears inside a naively written executor, and is removed
by scheduling alone.

\subsection{Circuit quality on a real workload}

Common subexpression elimination is a textbook pass, but its absence from the
path that feeds a backend is easy to miss: in our system it existed only
behind an explicit optimizer entry point, so every circuit handed to the FHE
and ZKP backends carried whatever redundancy the source contained. A fold
written the obvious way repeats its accumulator subtree at each step and
flattens \emph{exponentially}: $4372$ operations at $n{=}8$, collapsing to
$28$. On a real trained image classifier the same pass takes the emitted
circuit from $2365$ to $1764$ gates ($-25\%$) with bit-identical outputs.

\subsection{The same program, a different circuit per
backend}\label{sec:backend-aware}

Divergence on its own proves nothing: two circuits can differ in ways nobody
pays for. We therefore optimize one program once per backend and price
\emph{every} resulting circuit under \emph{every} backend
(\Cref{tab:crossapply}).

\begin{table}[t]
\centering
\caption{Cross-application, \bt{max-fold}, $n{=}8$, WAN, three parties.
Predicted milliseconds; $*$ marks each circuit under the backend it was
optimized for. Every circuit was differentially verified.}
\label{tab:crossapply}
\begin{tabular}{@{}lrrrl@{}}
\toprule
circuit optimized for & as \bt{shamir} & as \bt{semi-bin} & as \bt{bmr} & rules kept \\
\midrule
\bt{shamir} & $1453^{*}$ & $3761$ & $282$ & reassociate-reduction \\
\bt{semi-bin} & $1453$ & $3761^{*}$ & $282$ & reassociate-reduction \\
\bt{bmr} & $3253$ & $8761$ & $282^{*}$ & --- \\
\bottomrule
\end{tabular}
\end{table}

The diagonal wins every column. The mechanism is visible in the last column:
the \emph{same} rewrite is kept for the two interactive protocols and declined
for the constant-round garbled one, because depth is the round count for the
former and free for the latter. Running the garbled backend's circuit under
\bt{shamir} costs $2.24\times$. Nothing told the search this; the cost model
priced it.

Across three workloads --- an independent layer, a dependent reduction, and a
dot product against public constants --- one program yields $2$, $2$ and $3$
distinct circuits, with worst wrong-circuit penalties of $1.66\times$,
$2.33\times$ and $1.81\times$. Where a backend keeps no rules it genuinely
wanted the original circuit; an optimizer that always rewrites is not being
selective.

\subsection{A negative result: FHE does not host the layer}\label{sec:assign}

The same cost model assigns operation \emph{kinds} to protocols, and because
it spans FHE as well as MPC it can ask a question the MPC-only assignment
literature cannot: should the linear layer be an FHE program at all? Earlier
drafts of this work answered yes, reproducing the architecture of hand-designed
hybrid inference systems. Executing that answer showed it was wrong, and the
reason is worth more than the original claim.

\paragraph{Feasibility, not cost.} Switching a \emph{secret-shared} value into
FHE does not hand the engine the value. It hands it the un-reduced sum of the
parties' share components, and those components are uniform over the entire
ring --- a sharing of a small value over $\mathbb{Z}_{2^{64}}$ still has $64$-bit
shares. A ciphertext$\times$ciphertext product of two such values is congruent
to $a\cdot b$ modulo the ring, which is why the construction is correct, but as
an integer it is about $(\text{parties}\cdot 2^{k})^2$ for a $k$-bit ring. The
plaintext modulus must hold it: $2^{130}$ at a $64$-bit ring, against the
$2^{32}$ our depth-one parameters provide.

The model had a \emph{depth} budget but no \emph{modulus} budget, so it could
rank an assignment it was unable to perform. It now carries both, and refuses.
The gate's threshold agrees with the executed envelope --- the packed switch
works at $10$-bit shares and fails at $12$ --- without having been fitted to it.

\paragraph{And where it is feasible, it is not cheaper.} At a narrow ring
where the crossing can be performed, FHE is priced and simply not selected.
Executing the matched comparison at that width (two parties, $256$-wide layer,
\bt{semi2k} against the FHE route) agrees: MPC is faster and moves
$5.8\times$ fewer bytes.

\paragraph{What survives.} The comparison still moves to a binary sharing ---
the half of the classical hybrid split that concerns nonlinearities --- and
that decision is still deployment-dependent. What does not survive is the
linear half.

\paragraph{Why this does not contradict Gazelle and its successors.} Those
systems keep the model in \emph{plaintext}, so their linear layer is
ciphertext$\times$plaintext: no share magnitude is squared and the modulus
stays small. They also operate in a two-party client/server setting rather
than assigning operations within a multiparty computation, and in the
multiparty setting a public-constant multiply is already local, so there is
nothing for FHE to win. Our model represents the multiparty setting, and
within it the answer is negative. Hiding the \emph{model} as well as the input
--- the setting that makes both-operands-secret interesting --- is precisely what
puts the linear layer beyond a depth-one FHE engine.

\paragraph{Methodological note.} Getting here took four comparisons that
appeared to show the model wrong and instead showed two configurations being
compared that were not the same: an \bt{aby3} measurement against a
\bt{semi2k} prediction, and a $10$-bit-ring measurement against a $64$-bit-ring
prediction. A cost model spanning several backends has a large configuration
surface --- protocol, party count, threat model, network profile, value width,
ring width --- and any one of them silently makes two numbers incomparable while
both look reasonable. We added ring width to the model's vocabulary for this
reason: the one configuration the FHE route could actually be executed in was
one the model could not name.

\section{Related work}\label{sec:related}

\paragraph{Protocol assignment.} HyCC~\cite{HyCC}, Silph~\cite{Silph} and SING~\cite{SING} assign MPC protocols to
parts of a program and pay for the conversions. Silph solves the assignment
with integer linear programming~\cite{Silph}, which forces a linear cost model; SING~\cite{SING}
replaces the solver with a graph neural network, reporting large speedups at
comparable quality and supporting non-linear costs. Our objective is
non-separable for two structural reasons --- a dependence level costs the
\emph{maximum} over its members, and conversions batch per level --- so an ILP
formulation does not apply, and the exact model is cheap enough to evaluate
that a learned surrogate is not yet needed. We assign per operation
\emph{kind} rather than per instruction: a coarser space, but an interpretable
answer. \emph{These systems are MPC-only.}

\paragraph{Circuit compiler infrastructure.} CirC~\cite{CirC} shares our
one-IR-many-backends structure, with frontends and backends over a typed
term IR, and includes oblivious and linear-scan array elimination. Its
independent arrival at the same architecture is reassuring. CirC optimizes for
a target but does not, to our knowledge, report circuits that \emph{diverge}
per backend under a shared cost model.

\paragraph{Hand-built private inference.} Gazelle~\cite{Gazelle},
Delphi~\cite{Delphi}, Cheetah~\cite{Cheetah} and CrypTFlow2~\cite{CrypTFlow2}
are hand-designed systems that split linear layers and
nonlinearities across FHE and MPC. They are the expert artefacts this paper
argues a compiler should be able to derive; our assignment reproduces the
shape of that split from a cost model, though see \Cref{sec:limitations}.

\paragraph{MPC compilers.} MP-SPDZ~\cite{MPSPDZ} optimizes rounds by merging independent
openings --- the same insight our scheduler implements --- and is blind to
FHE. FHE compilers optimize multiplicative depth and are blind to MPC. Neither
can ask whether a computation should be an FHE program at all, because neither
has a cost model spanning the alternatives.

\section{Limitations}\label{sec:limitations}

\paragraph{Measured versus modelled.} \Cref{sec:eval-network} and
\Cref{sec:ablation} are executed: real serialization, real message passing,
real byte counting, verified against a plaintext reference.
\Cref{tab:crossapply} and the assignment results are \emph{predictions} from
the analytic cost model. The model's online round and volume terms are exact
for the lowering; its offline volume and local-work constants are literature
estimates, and we do not quote them as measurements.

\paragraph{Metric.} ``Interactive invocations'' counts batched protocol calls
at the backend interface, which is exactly comparable between routes but is a
proxy for rounds rather than a round count: one batched comparison runs
several rounds internally.

\paragraph{Scope of the compiled path.} The compiled network is MPC. Our
FHE/MPC hybrid remains orchestration outside the language --- the language does
not express scheme switching --- so the expert baseline in
\Cref{sec:eval-network} is the batched gadget library rather than the more
heavily tuned hybrid pipeline. Extending the language to express switching is
future work and would strengthen the backend-aware claim, since the compiler
could then choose where to switch.

\paragraph{Workload scale.} The compiled networks here are small MLPs.
Convolutional models at image scale are compiled and costed but not executed
end-to-end through the compiled path. \TODO{scale the compiled path to the
convolutional model, or state the ceiling analytically}

\paragraph{Optimality.} The rewrite search is greedy and site-independent; a
pair of rewrites that only pays jointly is missed. Assignment is greedy over
operation kinds and optimal assignment is NP-hard.

\section{Conclusion}

Private machine learning is written by hand against a chosen backend, and the
choice is usually made before the program is. We have argued that both
decisions should be compiled: that a schedule batches independent work more
reliably than an author does, and that the circuit worth emitting depends on
the backend it is headed for. The evidence for the first is that compiled code
beat our own experts on their own network, and that the defect it removes is
one those experts shipped three times in one codebase. The evidence for the
second is that a single rewrite is a $2.24\times$ win under one protocol and
declined under another, from one rule and one cost model.

\bibliographystyle{plain}
\bibliography{references}

\end{document}
